\documentclass[11pt,a4paper]{article}
\usepackage[T1]{fontenc}
\usepackage[utf8]{inputenc}
\usepackage{amsmath,amssymb,amsthm}
\usepackage[margin=1in]{geometry}
\usepackage[colorlinks=true,linkcolor=blue,urlcolor=blue]{hyperref}
\theoremstyle{plain}
\newtheorem{theorem}{Theorem}
\newtheorem{lemma}[theorem]{Lemma}
\newtheorem{proposition}[theorem]{Proposition}
\newtheorem{corollary}[theorem]{Corollary}
\theoremstyle{definition}
\newtheorem{remark}[theorem]{Remark}

\title{The empirical recurrence for OEIS A235966, proved by transfer matrix}
\author{Adrian Perez Fontelles\\ \small Independent researcher}
\date{2 September 2026}
\begin{document}
\maketitle

\begin{abstract}
OEIS A235966 counts arrays over $\{0,\dots,3\}$ subject to a condition on the order
statistics of every $2\times2$ subblock, and carries an empirical recurrence of order
$29$ contributed by R.~H.~Hardin. It is true, and it is decidable rather than empirical.
The condition is local to two consecutive rows, so the count is a walk count on $S=1964$
states and is $C$-finite; whether the conjectured recurrence annihilates it is then settled by
a finite exact computation, with the Cayley--Hamilton theorem bounding the work.
\end{abstract}

\noindent\small 2020 Mathematics Subject Classification. 05A15, 05B45, 62G30.\normalsize

\section{The conjecture}

OEIS A235966 is ``Number of (n+1)X(2+1) 0..3 arrays colored with the upper median value of each 2X2 subblock.''. It has offset $1$ and begins
\[
1964, 16744, 119770, 856720, 6207010, 44922684, 325144432, 2361022844,\ \dots
\]
The entry states, as a conjecture contributed by its author and never marked settled:
\begin{quote}
Empirical: a(n) = 160*a(n-3) +1405*a(n-4) +2240*a(n-5) +5011*a(n-6) -47892*a(n-7) -167088*a(n-8) -266910*a(n-9) -62360*a(n-10) +1453816*a(n-11) +3423362*a(n-12) +2700096*a(n-13) -2875681*a(n-14) -1353448*a(n-15) +2584789*a(n-16) -4251920*a(n-17) -8824295*a(n-18) +685061*a(n-19) +4397856*a(n-20) +158216*a(n-21) +4986700*a(n-22) +129756*a(n-23) +547072*a(n-24) +96320*a(n-25) -215488*a(n-26) -16688*a(n-27) -35392*a(n-28) -9216*a(n-29) for n>34
\end{quote}
As of the ``Last modified'' line on the live entry (Jul 23 2025 08:57:55, revision 6) this is still
recorded as empirical, and nothing on the entry records it as proved.

\section{What is being counted}

The $2\times2$ subblocks of an $(n+1)\times3$ array form an
$n\times2$ grid. Colour each of them by $\sigma=c$, where $a\le b\le c\le d$ are its four entries sorted --- its minimum, lower median, upper median and maximum. ``Coloured with'' is the entry's way of
asking for a PROPER colouring of that grid: subblocks that are neighbours, horizontally or
vertically, must receive different colours.

That reading is fixed by the entries themselves rather than guessed. Take the upper median
over $0..2$. A $2\times2$ array has one subblock and no neighbouring pair, so every one of the
$3^4=81$ arrays qualifies; a $2\times3$ array has one horizontal pair and $294$ of its $729$
arrays qualify; a $3\times3$ array has two horizontal and two vertical pairs and $722$
qualify. Those are the three numbers the corresponding entries publish. Imposing no condition
would give $729$ at $2\times3$, and counting the two diagonal pairs as neighbours as well
would give $0$ at $3\times3$, so neither is what is meant.

\section{The count is a walk count}

A subblock's colour is decided by two consecutive rows of the array, and
the vertical condition compares two consecutive rows of the subblock grid, so the state is the
PAIR of consecutive array rows. A step appends a row: that fixes a new subblock row, which
must be properly coloured along itself and must differ from the previous subblock row in every
position. An $(n+1)$-row array is a walk of $n-1$ steps, so
\[
a(n)\;=\;\iota^{\!\top}M^{\,n-1}\tau,\qquad\iota=\tau=(1,\dots,1)^{\!\top},
\]
over the $S=1964$ pairs whose own subblock row is properly coloured.
The alignment of the walk index with the entry's own $n$ was checked against its published
terms rather than assumed. In particular $a$ is $C$-finite.

\section{The proof}

\begin{lemma}\label{lem:crit}
Let $q(t)=t^{29}-\sum_i c_it^{29-i}$ be the characteristic polynomial of the
conjectured recurrence. The residual $a(n)-\sum_ic_ia(n-i)$ equals
$\iota^{\!\top}M^{\,j}q(M)\tau$ for the corresponding walk index $j$, so the recurrence
holds from some point on if and only if $\iota^{\!\top}M^{j}q(M)\tau=0$ for all large $j$,
and it is enough to examine $j<S$.
\end{lemma}

\begin{proof}
Factor $M^{j}$ out of $M^{j+29}-\sum_ic_iM^{j+29-i}$. For the bound, $M^{S}$
is an integer combination of $I,M,\dots,M^{S-1}$ by Cayley--Hamilton, so the numbers
$u_j=\iota^{\!\top}M^{j}q(M)\tau$ obey the monic recurrence given by the characteristic
polynomial of $M$; $S$ consecutive zeros force every later one, and the last nonzero $u_j$
pins the threshold exactly.
\end{proof}

\begin{theorem}
\[
a(n)\;=\;160\,a(n-3) + 1405\,a(n-4) + 2240\,a(n-5) + 5011\,a(n-6) - 47892\,a(n-7) - 167088\,a(n-8) - 266910\,a(n-9) - 62360\,a(n-10) + 1453816\,a(n-11) + 3423362\,a(n-12) + 2700096\,a(n-13) - 2875681\,a(n-14) - 1353448\,a(n-15) + 2584789\,a(n-16) - 4251920\,a(n-17) - 8824295\,a(n-18) + 685061\,a(n-19) + 4397856\,a(n-20) + 158216\,a(n-21) + 4986700\,a(n-22) + 129756\,a(n-23) + 547072\,a(n-24) + 96320\,a(n-25) - 215488\,a(n-26) - 16688\,a(n-27) - 35392\,a(n-28) - 9216\,a(n-29)
\]
for every $n>34$.
\end{theorem}

\begin{proof}
The vector $w=q(M)\tau$ was formed by $29$ matrix--vector products in exact integer
arithmetic and $\iota^{\!\top}M^{j}w$ evaluated for $j=0,1,\dots$, again exactly; those
integers vanish from the index corresponding to $n=34+1$ onwards and the run continued
until the working vector was identically zero, which settles every larger $j$ at once.
\end{proof}

\section{Verification}

Three checks, each independent of the algebra above.

First, the model reproduces all $17$ terms the entry publishes, exactly, in integer
arithmetic. This is what ties the model to the entry: the digraph is built by reading the
entry's English, and a misreading gives different counts.

Second, the conjectured recurrence was evaluated directly on those published terms, with no
matrices involved, and holds wherever the proved range and the data overlap.

Third, the annihilation test of Lemma~\ref{lem:crit} was carried out in exact integer
arithmetic on the whole vector rather than on a sampled prefix.

\begin{thebibliography}{9}
\bibitem{oeis} The OEIS Foundation, \emph{The On-Line Encyclopedia of Integer
Sequences}, \url{https://oeis.org}, 2026. Sequence A235966.
\bibitem{stanley} R.~P.~Stanley, \emph{Enumerative Combinatorics}, Volume 1, 2nd ed.,
Cambridge University Press, 2011. (Transfer-matrix method, Section 4.7.)
\bibitem{david} H.~A.~David and H.~N.~Nagaraja, \emph{Order Statistics}, 3rd ed., Wiley,
2003.
\bibitem{horn} R.~A.~Horn and C.~R.~Johnson, \emph{Matrix Analysis}, 2nd ed., Cambridge
University Press, 2013.
\end{thebibliography}

\end{document}
